\documentclass[a4paper]{article}
\usepackage{amsfonts}
\usepackage{amsmath}
\usepackage{amssymb}
\usepackage{graphicx}     

\begin{document}
  
\noindent \large Auteur: Anna Voronovska \\
\noindent \large Studierichting: Informatica\\
\noindent \large Oefeningenreeks 4A nr. 44.7\\

\medskip

\normalsize

$\displaystyle\int \dfrac{x^2-43}{x^3+x+10}dx$\\

Noemer ontbinden, $x=-2$:\\

$(-2)^3 + (-2) + 10 = -8 - 2 + 10 = 0$\\

Dus $(x+2)$ is een factor:\\

Horner deling met $x=-2$:\\

$\begin{array}{c|cccc}
   & 1 & 0  & 1 & 10 \\
-2 &   & -2 & 4 & -10 \\
\hline
   & 1 & -2 & 5 & 0
\end{array}$\\

Dus $x^3 + x + 10 = (x+2)(x^2-2x+5)$\\

$x^3 + x + 10 = (x+2)(x^2-2x+5)$\\

Partieelbreuken:\\

$\dfrac{x^2-43}{(x+2)(x^2-2x+5)} = \dfrac{A}{x+2} + \dfrac{Bx+C}{x^2-2x+5}$\\

$x^2-43 = A(x^2-2x+5) + (Bx+C)(x+2)$\\

$A = \left. \dfrac{x^2-43}{x^2-2x+5}\right|_{x=-2} = \dfrac{4-43}{4+4+5} = \dfrac{-39}{13} = -3$\\

$B(1+2i)+C \Rightarrow \left. \dfrac{x^2-43}{x+2}\right|_{x=1+2i} = \dfrac{(1+2i)^2-43}{(1+2i)+2} = \dfrac{1+4i-4-43}{3+2i} = \dfrac{-46+4i}{3+2i}$\\

$= \dfrac{(-46+4i)(3-2i)}{(3+2i)(3-2i)} = \dfrac{-138+92i+12i+8}{13} = \dfrac{-130+104i}{13} = -10+8i$\\

$B(1+2i)+C = -10+8i$\\

$\Rightarrow 2B = 8$\\

$\Rightarrow B = 4$\\

$B+C = -10$\\

$ \Rightarrow C = -14$\\

Dus:\\

$\dfrac{x^2-43}{(x+2)(x^2-2x+5)} = \dfrac{-3}{x+2} + \dfrac{4x-14}{x^2-2x+5}$\\

Integreren:\\

Eerste integraal:\\

$\displaystyle\int \dfrac{-3}{x+2}dx = -3\ln|x+2|$\\

Tweede integraal:\\

$\displaystyle\int \dfrac{4x-14}{x^2-2x+5}dx = \displaystyle\int \dfrac{4x-4}{x^2-2x+5}dx - \displaystyle\int \dfrac{10}{(x-1)^2+4}dx$\\

Eerste deel: $u = x^2-2x+5 \Rightarrow du = (2x-2)dx$:\\

$\displaystyle\int \dfrac{4x-4}{x^2-2x+5}dx = 2\displaystyle\int \dfrac{du}{u} = 2\ln|x^2-2x+5|$\\

Tweede deel: $t = x-1$:\\

$\displaystyle\int \dfrac{10}{(x-1)^2+4}dx = \dfrac{10}{2}\operatorname{Bgtan}\left(\dfrac{x-1}{2}\right) = 5\operatorname{Bgtan}\left(\dfrac{x-1}{2}\right)$\\

Alles samen:\\

$\displaystyle\int \dfrac{x^2-43}{x^3+x+10}dx = -3\ln|x+2| + 2\ln(x^2-2x+5) - 5\operatorname{Bgtan}\left(\dfrac{x-1}{2}\right) + c$




\begin{thebibliography}{999}
%\bibitem{a} Referentie
\end{thebibliography}
\end{document} 
